\documentclass[11pt]{article}

\usepackage{amsmath,amssymb,amsthm}

\newtheorem{theorem}{Theorem}
\newtheorem{lemma}[theorem]{Lemma}
\newtheorem{corollary}[theorem]{Corollary}
\theoremstyle{definition}
\newtheorem{definition}[theorem]{Definition}
\newtheorem{remark}[theorem]{Remark}

\title{Disproof of conjecture \texttt{00000000427}}
\author{Submission \texttt{earthking11\_submission\_20260912150204}}
\date{12 September 2026}

\begin{document}

\maketitle

\begin{abstract}
Conjecture \texttt{00000000427} asserts that the count of spin symmetric plane
partitions at $t=1$ is given by
$F(n) = 2^{\lfloor n^2/4\rfloor}\prod_{i=1}^{n}(2i-1)!!/i!$. We show that
$F(n)$ is a non-integer rational number for every $n\ge 4$; in particular
$F(4)=525/2$. Since a count at $t=1$ is a cardinality, and hence a
non-negative integer, the conjecture is false already at $n=4$. An independent
exhaustive enumeration of order ideals of $[n]^3$ is included as supporting
evidence, and does not enter the proof.
\end{abstract}

\section{The conjecture}

\begin{quote}
\textbf{Definition.} Spin symmetric plane partitions are plane partitions in a
cubic box symmetric under all diagonals and weighted by spin.\\[2pt]
\textbf{Conjecture.} Their count at $t = 1$ equals
\[
  2^{\lfloor n^{2}/4\rfloor}\prod_{i=1}^{n}\frac{(2i-1)!!}{i!},
\]
and the ratio of this formula to the unsymmetrized MacMahon formula is exactly a
product of a power of two and double factorials, with no other prime factors.
\end{quote}

The conjectured formula, read as an exact rational number, is the object we
analyse.

\section{The formula is not an integer}

\begin{definition}\label{def:F}
The \emph{odd double factorial} is
$(2i-1)!! = \prod_{j=1}^{i}(2j-1)$ for $i\ge 1$, with $(2\cdot 0-1)!! = 1$.
For $n\ge 0$ define
\[
  F(n) \;=\; 2^{\lfloor n^{2}/4\rfloor}\prod_{i=1}^{n}\frac{(2i-1)!!}{i!}
  \;\in\; \mathbb{Q}.
\]
Write $\nu_{2}$ for the $2$-adic valuation of a non-zero rational, and set
$D(n) = \sum_{i=1}^{n}\nu_{2}(i!)$.
\end{definition}

\begin{lemma}\label{lem:valuation}
For every $n\ge 0$,
\[
  \nu_{2}\bigl(F(n)\bigr) \;=\; \Bigl\lfloor \frac{n^{2}}{4}\Bigr\rfloor - D(n).
\]
\end{lemma}

\begin{proof}
Every $(2i-1)!!$ is a product of odd integers, hence odd, so
$\nu_{2}\bigl((2i-1)!!\bigr) = 0$. The valuation is additive on products and
satisfies $\nu_{2}(x/y) = \nu_{2}(x)-\nu_{2}(y)$ for non-zero rationals
$x,y$. Therefore
\[
  \nu_{2}\Bigl(\prod_{i=1}^{n}\frac{(2i-1)!!}{i!}\Bigr)
  \;=\; \sum_{i=1}^{n}\bigl(0-\nu_{2}(i!)\bigr) \;=\; -D(n),
\]
and, since the leading factor $2^{\lfloor n^{2}/4\rfloor}$ is a power of two,
\[
  \nu_{2}\bigl(F(n)\bigr)
  \;=\; \Bigl\lfloor \frac{n^{2}}{4}\Bigr\rfloor + \bigl(-D(n)\bigr)
  \;=\; \Bigl\lfloor \frac{n^{2}}{4}\Bigr\rfloor - D(n).
\]
This is the claim.
\end{proof}

\begin{theorem}[The counterexample $n=4$]\label{thm:four}
$F(4) = 525/2$, which is not an integer. In particular
$\nu_{2}\bigl(F(4)\bigr) = -1$.
\end{theorem}

\begin{proof}
We compute
\[
  D(4) \;=\; \nu_{2}(1!)+\nu_{2}(2!)+\nu_{2}(3!)+\nu_{2}(4!)
  \;=\; 0+1+1+3 \;=\; 5,
\]
while $\lfloor 4^{2}/4\rfloor = 4$. By Lemma~\ref{lem:valuation},
\[
  \nu_{2}\bigl(F(4)\bigr) \;=\; 4 - 5 \;=\; -1 < 0,
\]
so $F(4)$ has a factor of $2$ in its denominator in lowest terms, hence is not
an integer. Directly,
\[
  F(4) \;=\; 2^{4}\cdot\frac{1\cdot 3\cdot 15\cdot 105}{1!\cdot 2!\cdot 3!\cdot 4!}
  \;=\; 16\cdot\frac{4725}{288}
  \;=\; \frac{75600}{288}
  \;=\; \frac{525}{2}. \qedhere
\]
\end{proof}

\begin{theorem}\label{thm:all}
For every $n\ge 4$, $F(n)$ is not an integer; more precisely
$\nu_{2}\bigl(F(n)\bigr) < 0$.
\end{theorem}

\begin{proof}
Put $G(n) = \lfloor n^{2}/4\rfloor - D(n) = \nu_{2}(F(n))$ by
Lemma~\ref{lem:valuation}. We first compute the increment. Since
$n^{2}-(n-1)^{2} = 2n-1$ one checks
\[
  \Bigl\lfloor \frac{n^{2}}{4}\Bigr\rfloor - \Bigl\lfloor \frac{(n-1)^{2}}{4}\Bigr\rfloor
  \;=\; \Bigl\lfloor \frac{n}{2}\Bigr\rfloor,
\]
and $D(n)-D(n-1) = \nu_{2}(n!)$. Legendre's formula gives
$\nu_{2}(n!) = \sum_{j\ge 1}\lfloor n/2^{j}\rfloor$. Hence
\[
  G(n)-G(n-1)
  \;=\; \Bigl\lfloor \frac{n}{2}\Bigr\rfloor - \nu_{2}(n!)
  \;=\; -\sum_{j\ge 2}\Bigl\lfloor \frac{n}{2^{j}}\Bigr\rfloor \;\le\; 0 .
\]
For $n\ge 4$ we have $\lfloor n/4\rfloor \ge 1$, so $G(n) < G(n-1)$: the
sequence $G$ is strictly decreasing from $n=4$ onward. Finally
$G(3) = \lfloor 9/4\rfloor - D(3) = 2 - (0+1+1) = 0$, so $G(4) = -1$ and by
induction $G(n)\le -1$ for all $n\ge 4$. Thus
$\nu_{2}(F(n)) = G(n) < 0$ for every $n\ge 4$.
\end{proof}

\begin{corollary}\label{cor:false}
No natural number equals $F(4)=525/2$. More generally, the conjectured
equality between $F(n)$ and the count of objects at $t=1$ fails for every
$n\ge 4$, and in particular at $n=4$. Conjecture \texttt{00000000427} is
therefore false.
\end{corollary}

\begin{proof}
A count at $t=1$ is the cardinality of a finite set, hence a non-negative
integer. But $F(4)=525/2\notin\mathbb{N}$ by Theorem~\ref{thm:four}, and
$F(n)\notin\mathbb{Z}$ for all $n\ge4$ by Theorem~\ref{thm:all}. Hence no such
count can equal the conjectured formula.
\end{proof}

\begin{remark}
For $n = 1,2,3$ the formula is integral,
$F(1)=1$, $F(2)=3$, $F(3)=15$, which is presumably why the claim looked
plausible; integrality fails at $n=4$ and, by Theorem~\ref{thm:all}, never
recovers. The valuation mechanism is that $\nu_{2}(n!) = n - s_{2}(n)$ (where
$s_{2}(n)$ is the number of ones in the binary expansion of $n$) grows faster
than $\lfloor n/2\rfloor$, so the exponent of $2$ in the denominator of $F(n)$
strictly outruns the explicit power $2^{\lfloor n^{2}/4\rfloor}$.
\end{remark}

\begin{remark}
The secondary claim about the ratio to MacMahon's formula also fails under the
natural reading of ``a product of a power of two and double factorials'' as a
non-negative integer: at $n=4$, using the MacMahon total
$\prod_{i,j,k=1}^{4}\frac{i+j+k-1}{i+j+k-2} = 232848$, the ratio is
\[
  \frac{F(4)}{232848} \;=\; \frac{525}{2\cdot 232848}
  \;=\; \frac{525}{465696},
\]
whose numerator is odd while its denominator is even. The main refutation,
however, rests only on the count equality.
\end{remark}

\section{Independent enumeration (supporting evidence)}

As a cross-check independent of the valuation argument, we exhaustively
enumerated order ideals of the product poset $[n]^{3}$ (equivalently, plane
partitions in the $n$-cube) and counted those invariant under all permutations
of the three coordinates.

\begin{center}
\begin{tabular}{c|c|c}
$n$ & order ideals of $[n]^{3}$ & invariant under all coordinate permutations \\
\hline
1 & 2     & 2   \\
2 & 20    & 5   \\
3 & 980   & 16  \\
4 & 232848 & 66
\end{tabular}
\end{center}

The MacMahon totals $2,20,980,232848$ reproduce the classical product formula
and validate the enumerator. The invariant counts $2,5,16,66$ disagree with the
conjectured formula $1,3,15,525/2$ already at $n=1$ ($1$ vs $2$) and $n=2$
($3$ vs $5$), and at $n=4$ the conjectured value is not even an integer.

The exact definition of ``spin symmetric plane partition'' is not pinned down
by the conjecture, so this enumeration is offered as supporting evidence and
not as the definition. The refutation does not rely on it: whatever the
intended object, its $t=1$ value is a count, and no count equals $525/2$.

\section{Conclusion}

The conjectured formula $F(n)$ is not an integer for any $n\ge 4$: already
$F(4)=525/2$, and $\nu_{2}(F(n))<0$ for all $n\ge 4$. A count at $t=1$ is a
non-negative integer, so the conjectured equality is false. Conjecture
\texttt{00000000427} is disproved.

\end{document}
